\begin{pastexam}[label={px:gw3-096-ds-05}]{關務三等資訊處理 96 年 第五題}
（申論；25 分）

下列為一個環狀串列（Circular Queue）中加入與刪除一個元素的演算法：
\begin{quote}
\begin{tabbing}
\hspace*{1.5em}\=\kill
Procedure ADDQ(item,Q,n,front,rear)\\
\> rear$\leftarrow$(rear+1) mod n\\
\> if front=rear then call QUEUE\_FULL\\
\> Q(rear)$\leftarrow$item\\
end ADDQ\\
Procedure DELETEQ(item,Q,n,front,rear)\\
\> if front=rear then call QUEUE\_EMPTY\\
\> front$\leftarrow$(front+1) mod n\\
\> item$\leftarrow$Q(front)\\
end DELETEQ
\end{tabbing}
\end{quote}
\begin{itemize}[leftmargin=2.2em, nosep]
\item[㈠] 在演算法 ADDQ 中當 front=rear 時發生溢位（Overflow），實際上還有一個空間，請說明為何不使用呢？（5 分）
\item[㈡] 若欲使用此一空間，則 ADDQ 及 DELETEQ 應如何改寫，試寫出他們修改後的 ADDQ 及 DELETEQ 演算法。（20 分）
\end{itemize}
\end{pastexam}

\begin{pastsoln}
答案：㈠ 因為空佇列的條件就是 $\mathit{front} = \mathit{rear}$，把最後一格也填滿會使滿佇列同樣是 $\mathit{front} = \mathit{rear}$，
分不出空還是滿。 ㈡ 多用一個旗標 $\mathit{tag}$ 記錄「上一次讓兩者相等的是加入還是刪除」，演算法在下面。

\textbf{題目演算法的約定。} $Q(0..n-1)$，$\mathit{front}$ 指向第一個元素的\textbf{前一格}，$\mathit{rear}$ 指向最後一個
元素，開始時 $\mathit{front} = \mathit{rear} = 0$。 刪除時先把 $\mathit{front}$ 前進一格再取值。 元素個數是
$(\mathit{rear} - \mathit{front}) \bmod n$（資結書第~6.2~節）。

\textbf{㈠} 空佇列的判斷是 $\mathit{front} = \mathit{rear}$。 假設允許放滿 $n$ 個：從空佇列（$\mathit{front} = \mathit{rear}$）
連續加入 $n$ 個，$\mathit{rear}$ 前進 $n$ 格繞了一整圈，又回到 $\mathit{front}$，滿的時候也是 $\mathit{front} = \mathit{rear}$。
兩個指標只有 $n$ 種相對位置（差 $0, 1, \dots, n - 1$），卻要表示 $0$ 到 $n$ 共 $n + 1$ 種元素個數，一定有兩種撞在一起。
所以 ADDQ 在 $\mathit{rear}$ 前進後碰到 $\mathit{front}$ 就宣告溢位，犧牲 $\mathit{front}$ 那一格，最多只放 $n - 1$ 個，
空和滿才分得開。 （順帶一提，原來的 ADDQ 在宣告溢位前已經把 $\mathit{rear}$ 往前移了，QUEUE\_FULL 處理時應該把它退回
原位，否則佇列狀態被破壞；下面的寫法先判斷再移動，沒有這個問題。）

\textbf{㈡ 改寫。} 加一個變數 $\mathit{tag}$：$\mathit{front} = \mathit{rear}$ 時，$\mathit{tag} = 1$ 表示滿、$\mathit{tag} = 0$ 表示空。
開始時 $\mathit{front} = \mathit{rear} = 0$、$\mathit{tag} = 0$。
\begin{pseudo}{使用全部 $n$ 格的環狀佇列}
\Procedure{AddQ}{$\mathit{item}, Q, n, \mathit{front}, \mathit{rear}, \mathit{tag}$}
    \If{$\mathit{front} = \mathit{rear}$ 且 $\mathit{tag} = 1$}
        \State \Call{Queue-Full}{} \Comment{$n$ 格都滿了}
    \EndIf
    \State $\mathit{rear} \gets (\mathit{rear} + 1) \bmod n$
    \State $Q(\mathit{rear}) \gets \mathit{item}$
    \If{$\mathit{front} = \mathit{rear}$}
        \State $\mathit{tag} \gets 1$ \Comment{加入後追上 front：滿}
    \EndIf
\EndProcedure
\Statex
\Procedure{DeleteQ}{$\mathit{item}, Q, n, \mathit{front}, \mathit{rear}, \mathit{tag}$}
    \If{$\mathit{front} = \mathit{rear}$ 且 $\mathit{tag} = 0$}
        \State \Call{Queue-Empty}{}
    \EndIf
    \State $\mathit{front} \gets (\mathit{front} + 1) \bmod n$
    \State $\mathit{item} \gets Q(\mathit{front})$
    \State $\mathit{tag} \gets 0$ \Comment{刪除後一定不滿}
\EndProcedure
\end{pseudo}
（QUEUE\_FULL、QUEUE\_EMPTY 和原題一樣，處理錯誤後不再往下執行。）

\textbf{正確性。} 維持的性質是：$\mathit{front} = \mathit{rear}$ 時，$\mathit{tag} = 1$ 恰好表示佇列有 $n$ 個元素，
$\mathit{tag} = 0$ 恰好表示 $0$ 個；$\mathit{front} \ne \mathit{rear}$ 時 $\mathit{tag} = 0$。
\begin{itemize}
\item 開始時空佇列，$\mathit{tag} = 0$，成立。
\item AddQ 只在不滿時執行。 加入後元素個數變成 $k + 1$；$\mathit{rear}$ 追上 $\mathit{front}$ 正好表示 $k + 1 = n$，此時設
$\mathit{tag} = 1$；否則 $\mathit{front} \ne \mathit{rear}$，$\mathit{tag}$ 原本是 $0$（不滿），維持 $0$。
\item DeleteQ 只在不空時執行。 刪除後個數 $\le n - 1$，不可能滿，設 $\mathit{tag} = 0$；若此時 $\mathit{front} = \mathit{rear}$，
就表示剛好刪成空佇列，和 $\mathit{tag} = 0$ 的意思一致。
\end{itemize}
於是 $n$ 格都能用。 每個操作仍是 $O(1)$，代價是每次多一個判斷和指定。 另一種做法是用一個計數器 $\mathit{count}$ 記錄
元素個數，$\mathit{count} = 0$ 為空、$\mathit{count} = n$ 為滿，效果相同。
\end{pastsoln}
